\section*{Chapter \ref{ch:s1:short-term-reversal} --- Short-Term Reversal}

\begin{solution}{exo:s1:short-term-reversal:1}
The weight traded is twice the turnover, so the break-even cost is $0.303/(2 \times 0.76 \times 252) = 7.9$ basis points per unit traded. At ten basis
points the cost is about $2 \times 0.75 \times 0.001 \times 252 = 37.8\%$ a year, more than the gross return.
\end{solution}

\begin{solution}{exo:s1:short-term-reversal:2}
The recorded return is the spread, $0.001 = 10$ basis points, with no change in value. A reversal signal on trade prices sees a winner and sells it; the next
day's trade at the bid would ``confirm'' the reversal. That is bid--ask bounce: a reversal in the data that no one can trade, which is why signals are built on
midquotes or closing auction prices, and evaluated on the prices the book could actually trade at.
\end{solution}

\begin{solution}{exo:s1:short-term-reversal:3}
$0.7 \times 0.02 \times \sqrt{10/25} = 0.00885$: 88.5 basis points, on top of the half spread.
\end{solution}

\begin{solution}{exo:s1:short-term-reversal:4}
$0.04 \times 0.02 \times 2 = 0.0016$, 16.0 basis points for one day. A trade of a hundredth of daily volume, \$0.25 million, costs
$2 + 10^4 \times 0.7 \times 0.02 \times \sqrt{0.25/25} = 16.0$ basis points: that is where the forecast stops paying. The \$10 million trade of exercise~3 would
cost more than five times the edge.
\end{solution}

\begin{solution}{exo:s1:short-term-reversal:5}
The Sharpe ratio depends on risk as well as return. The raw book's losers and winners share industries and styles, so the book carries factor bets that add
volatility without adding return; the residual book is nearly free of them. The return rises a little (27.0\% to 30.3\%), the volatility falls a lot.
\end{solution}

\begin{solution}{exo:s1:short-term-reversal:6}
On an announcement day the stock's move is mostly news, and in a market with post-earnings drift the news keeps moving the price the same way for weeks: the
reversal book buys a stock that keeps falling. Announcement dates are scheduled and published in advance, and the filter uses only whether the day was an
announcement day, which is known at the close, not the surprise's future effect.
\end{solution}

\begin{solution}{exo:s1:short-term-reversal:7}
The net Sharpe ratio at \$300 million is 0.73. Beyond \$1 billion the net ratio is within noise of zero (0.09, $-0.08$, 0.19 at \$1, \$3 and \$10 billion):
the book trades so little that its net return is a few basis points a year, and eight years of data cannot rank such small numbers. What is robust is the
fall from 1.72 at \$100 million to about zero at \$1 billion.
\end{solution}

\begin{solution}{exo:s1:short-term-reversal:8}
The Sharpe ratio of 13 is before costs; after ten basis points the same book's is $-3.43$. The loss at size is not a halving: at \$1 billion even the
cost-aware book earns its costs and nothing more. The right questions are the break-even cost, the book's turnover, and the capacity curve at the firm's
fitted costs.
\end{solution}

\begin{solution}{pb:s1:short-term-reversal:1}
\begin{enumerate}
\item Recent losers outperform recent winners over the next short window; a trader who needs immediacy moves the price, and whoever absorbs the flow is paid when
it comes back.
\item Jegadeesh: significant negative first-order serial correlation in monthly returns, 2.49\% a month between extreme deciles over 1934--1987. Lehmann: weekly
winners and losers reverse the next week, with profits that survive spreads and plausible costs.
\item Reversal returns behave like the returns to liquidity provision: predictable with the VIX, with expected returns and Sharpe ratios spiking in turmoil.
\item News changes the value; only the part of a move that is pressure is expected to come back.
\item 0.037, 0.038 and 0.041; the truth 0.044 (the same in the 500 most liquid, 0.039 for the industry-adjusted).
\item 6.69, 7.24 and 12.96 before costs; $-2.62$, $-2.46$ and $-3.43$ after ten basis points.
\item It is free of the industry and style exposures that add risk without return.
\item \omterm{def:s1:short-term-reversal:residual}{Residual reversal} earns risk-adjusted returns twice as large as conventional reversal and survives costs in large caps after 1990; removing the part
explained by cash-flow news made the risk-adjusted return four times the standard one's.
\item Grinold's rule, $\alpha = \mathrm{IC}\,\sigma z$, with $z$ the residual scaled by its own volatility and standardised, the IC measured before the
evaluation.
\item Expected return less a diagonal risk penalty less the half spread and square-root impact, neutral to the fifteen factors; with diagonal risk each name's
trade is a closed-form proximal step, and the neutrality multipliers come from a small Newton iteration.
\item 186\% of capital a day at a gross exposure of 2.48; 246\%, 738\% and 2,293\% of capital a year in costs at the three sizes.
\item It skips names on their announcement day; at \$1 billion it turns $-0.70\%$ a year into $+0.05\%$.
\item \textbf{Named result.} Net Sharpe ratios of 1.72 at \$100 million and 0.09 at \$1 billion with costs inside the optimiser, against $-27.8$ and $-31.7$ for
the book that ignores them.
\item Below \$1 billion for this signal: the edge per trade is about sixteen basis points for a name two standard deviations out, and square-root impact reaches
that at 1\% of a name's daily volume.
\item It earned 10.9\% a year (Sharpe ratio 0.93) to 1989 and 1.6\% (0.13) from 1990, before costs; $-5.2\%$ a year in the 2020s to date.
\item It is a monthly, value-weighted factor on large portfolios; the daily residual effect is a different strategy, and reversal returns concentrate in crises.
\item Reversal profits are earned overnight, with continuation within each part of the day; intraday reversal is about imbalances of less than an hour and bid--ask
bounce.
\item With \$50 million the daily residual book; with \$5 billion the overlay on a slower book, which costs almost nothing of its own.
\item The cost model: the strategy's sign after costs depends on it.
\item Being paid a small fee, many times a day, for trading with people in a hurry, and keeping more of it than it costs to trade.
\end{enumerate}
\end{solution}

\begin{solution}{iq:s1:short-term-reversal:1}
Traders who need to trade now move prices beyond value; the price comes back as the imbalance is absorbed. The reversal trader is paid by those hurried traders,
as a daily-horizon market maker, and is paid most when liquidity is scarce.
\end{solution}

\begin{solution}{iq:s1:short-term-reversal:2}
Raw reversal uses the whole return; \omterm{def:s1:short-term-reversal:residual}{residual reversal} removes what the factor model explains. Common moves are mostly news and do not reverse, and they give the
raw book factor exposures that add risk; the residual book is cleaner and, on US data and here, has a much higher risk-adjusted return.
\end{solution}

\begin{solution}{iq:s1:short-term-reversal:3}
Not as a naive daily book: a break-even of 15 basis points is close to real costs for most names. Ask for the cost model and its fit on real fills, the turnover,
the capacity curve with costs inside the optimiser, performance by liquidity bucket and around news, and its behaviour in crises.
\end{solution}

\begin{solution}{iq:s1:short-term-reversal:4}
Use a factor or diagonal risk model so the problem separates or nearly separates by name; the linear and power costs have a closed-form proximal step; handle
neutrality with a few Lagrange multipliers found by Newton's method, or use a proximal-gradient method (FISTA) on the factor model. Warm start from yesterday's
solution.
\end{solution}

\begin{solution}{iq:s1:short-term-reversal:5}
In a crisis when the flows it absorbs are informed or persistent (forced selling that goes on for days), when news days leak into the signal, and when a crowd
of similar books unwinds. Monitor losses on news days, exposure to the factors, turnover and costs against the model, and the book's correlation with other
reversal books.
\end{solution}

\begin{solution}{iq:s1:short-term-reversal:6}
The cost per unit traded of a trade of $Q$ dollars is $h + \eta\sigma\sqrt{Q/V}$. Setting it equal to $\alpha$ gives $Q^* = V\left((\alpha - h)/(\eta\sigma)\right)^2$
for $\alpha > h$, and no trade otherwise. With $\alpha$ = 16 basis points, $h$ = 2, $\eta = 0.7$ and $\sigma$ = 2\%, $Q^* = V/100$.
\end{solution}
